% Bond Pricing, Duration & Convexity — research-note PDF.
%
% PROSE, CODE AND TABLES ARE VERBATIM from the website article:
%   site/src/content/articles/bond-duration-convexity.tsx
%   site/src/content/articles/bond-code.ts
% Metadata from lib/articles.ts; project-card fallbacks from lib/topics.ts.
%
% CHARTS generated by figures/make_bond_duration_convexity.py from the article's
% own data module — the computed results the site's <Line> and <Bar> render.
\documentclass{pyportfolios-note}

\noteshorttitle{Bond Duration \& Convexity}

\begin{document}
\thispagestyle{notecover}

\notemasthead{Quantitative Insights}{Tutorial}{Q3 2026}{}

\notetitle{Bond Pricing, Duration\\\& Convexity}

\notedeck{Price, duration and convexity — the three numbers every fixed-income
desk lives by, verified against the worst bond year in modern history.}

\notelede{A bond is just a promise of future cash flows — so its price is
entirely a question of \textit{discounting}. From that single idea follow the
three numbers every fixed-income desk lives by: \textbf{price},
\textbf{duration} (first-order rate sensitivity, the bond's Delta) and
\textbf{convexity} (second-order, the bond's Gamma).}

\notecard
  {Quant Foundations \& Derivatives}
  {NumPy · Pandas · Matplotlib}
  {SHY, IEF, TLT (US Treasury ETFs)}
  {2022 (rate-shock year)}
  {Measuring and stress-testing interest-rate risk in a bond book}

% ============================================================= 1 Summary =====
\noteneedroom{150bp}
\notesection{1.\notenumsep Summary}

\begin{multicols}{2}
A bond's price is the present value of its coupons and principal, discounted at
the \textbf{yield to maturity}. Price and yield move inversely, and the
relationship is \textit{curved}, not linear. \textbf{Duration} measures the
percentage price change for a small yield change; \textbf{convexity} corrects
that estimate for the curvature — and the correction always works in the
bondholder's favour.
\end{multicols}

% =========================================================== 2 Intuition =====
\noteneedroom{170bp}
\notesection{2.\notenumsep Intuition}

\begin{multicols}{2}
\notesubsection{2.1\notenumsep Discounting, not magic}
When market yields rise from 2\% to 4\%, an old bond paying 2\% coupons is
suddenly a bad deal — nobody pays face value for below-market cash flows. Its
price falls until a buyer earns the \textit{new} market yield. No sentiment
involved: pure arithmetic.

\notesubsection{2.2\notenumsep Duration = weighted waiting time}
Macaulay duration is the average time you wait for your money, weighting each
cash flow by its share of present value. A 30y bond makes you wait decades →
huge sensitivity. A 2y note returns your money quickly → barely reacts.
\textbf{Rule of thumb: a bond loses \textasciitilde duration × Δy percent} when
yields rise by Δy.

\notesubsection{2.3\notenumsep Convexity: the curvature that helps you}
The price-yield curve is convex: losses from rising yields are \textit{smaller}
than the linear estimate, gains from falling yields are \textit{larger}. Long
bonds have the most convexity — which is why duration alone increasingly
misleads for big yield moves.
\end{multicols}

% ================================================= 3 Theory & Mechanics ======
\noteneedroom{190bp}
\notesection{3.\notenumsep Theory \& Mechanics}

With face $F$, coupon rate $c$ (paid $f$ times a year), yield $y$ and
$n = fT$ periods:

\[ P = \sum_{i=1}^{n} \frac{cF/f}{(1+y/f)^i} + \frac{F}{(1+y/f)^n} \]

\textbf{Macaulay duration} (years):
$D_{mac} = \frac{1}{P}\sum_i t_i \cdot PV(CF_i)$, and \textbf{modified duration}
$D = D_{mac}/(1+y/f)$.

\textbf{Convexity} $C$ is the second derivative of price w.r.t. yield, scaled by
price. Both derivatives can also be computed numerically — which is how we'll
verify the analytic formulas:

\[ \frac{\Delta P}{P} \approx -D\,\Delta y + \tfrac{1}{2} C (\Delta y)^2 \]

\begin{notecode}
def bond_price(face, coupon, y, T, freq=2):
    """Price of a fixed-coupon bond (coupon = annual rate, y = YTM)."""
    n = int(round(T * freq))
    t = np.arange(1, n + 1)
    cf = np.full(n, face * coupon / freq)
    cf[-1] += face
    return float(np.sum(cf / (1 + y / freq) ** t))

def duration_convexity(face, coupon, y, T, freq=2, h=1e-4):
    """Modified duration (years) and convexity via central differences."""
    p0 = bond_price(face, coupon, y, T, freq)
    p_up, p_dn = bond_price(face, coupon, y + h, T, freq), bond_price(face, coupon, y - h, T, freq)
    dur  = -(p_up - p_dn) / (2 * h * p0)
    conv =  (p_up - 2 * p0 + p_dn) / (h**2 * p0)
    return dur, conv

# quick check against the rule of thumb
p = bond_price(100, 0.04, 0.04, 10)
d, c = duration_convexity(100, 0.04, 0.04, 10)
print(f"10y 4% bond at par: P = {p:.2f}, duration = {d:.2f}y, convexity = {c:.1f}")
\end{notecode}

% ================================= 4 Applied Example — US Treasuries =========
\noteneedroom{200bp}
\notesection{4.\notenumsep Applied Example — US Treasuries}

\notesubsection{4.1\notenumsep The price-yield curve, by maturity}

Same 4\% coupon, three maturities. Watch two things: the \textbf{slope} (longer =
steeper = more duration) and the \textbf{curvature} (longer = more convex).

\begin{notecode}
yields = np.linspace(0.001, 0.10, 200)
fig, ax = plt.subplots(figsize=(10, 5))
for T, color in [(2, "steelblue"), (10, "darkorange"), (30, "crimson")]:
    prices = [bond_price(100, 0.04, y, T) for y in yields]
    ax.plot(yields * 100, prices, color=color, label=f"{T}y Treasury")
ax.axhline(100, color="gray", ls=":", lw=1)
ax.axvline(4, color="gray", ls=":", lw=1)
ax.set_xlabel("Yield (%)"); ax.set_ylabel("Price")
ax.set_title("Price vs yield — 4% coupon, three maturities")
ax.legend()
plt.tight_layout(); plt.show()
\end{notecode}

\notefigure{figures/bond-duration-convexity/fig1-price-yield.png}
  {Figure 4.1 · Price vs yield — 4\% coupon, three maturities}
  {The price curve against the straight-line duration estimate; the gap between
   them is convexity}
  {Source: Author calculations, from the article's computed results.}

\notesubsection{4.2\notenumsep The duration ladder: a +100bp shock}

The whole risk story in one table — duration grows with maturity, and so does the
damage from the same yield move.

\begin{notecode}
rows = []
for T in [1, 2, 5, 10, 20, 30]:
    d, c = duration_convexity(100, 0.04, 0.04, T)
    p0 = bond_price(100, 0.04, 0.04, T)
    p1 = bond_price(100, 0.04, 0.05, T)          # +100bp, exact
    rows.append([T, round(d, 2), round(c, 1), f"{(p1/p0 - 1)*100:.2f}%"])

ladder = pd.DataFrame(rows, columns=["Maturity (y)", "Mod. duration", "Convexity", "Exact dP for +100bp"])
print(ladder.to_string(index=False))
\end{notecode}

\begin{notetable}{@{}p{120bp}p{120bp}p{110bp}p{146bp}@{}}
\thcell{Maturity (y)} & \thcell{Mod. duration} & \thcell{Convexity} & \thcell{Exact dP for +100bp} \\
\theadrule
1 & 0.97 & 1.4 & -0.96\% \\ \rowrule
2 & 1.90 & 4.6 & -1.88\% \\ \rowrule
5 & 4.49 & 23.5 & -4.38\% \\ \rowrule
10 & 8.18 & 78.9 & -7.79\% \\ \rowrule
20 & 13.68 & 239.9 & -12.55\% \\ \rowrule
30 & 17.38 & 420.8 & -15.45\% \\
\end{notetable}

\notesubsection{4.3\notenumsep Reality check: SHY, IEF, TLT in 2022}

2022 was the worst bond year in modern history: the 2y yield rose
\textasciitilde370bp, the 10y \textasciitilde237bp, the 30y
\textasciitilde207bp. If duration analytics mean anything, the three Treasury
ETFs' losses should line up with their durations (\textasciitilde1.9 /
\textasciitilde7.5 / \textasciitilde17.5). Let's check against real prices.

\begin{notecode}
def load_prices(ticker, start, end):
    """Adjusted-close prices: yfinance first, Stooq as fallback."""
    try:
        import yfinance as yf
        df = yf.download(ticker, start=start, end=end, auto_adjust=True, progress=False)
        if not df.empty:
            return df["Close"].squeeze().rename(ticker).dropna()
    except Exception as exc:
        print(f"yfinance failed ({exc}); trying Stooq…")
    url = f"https://stooq.com/q/d/l/?s={ticker.lower()}.us&i=d"
    df = pd.read_csv(url, parse_dates=["Date"], index_col="Date")
    return df.loc[start:end, "Close"].rename(ticker).dropna()

etfs = {"SHY": (1.9, 0.0370), "IEF": (7.5, 0.0237), "TLT": (17.5, 0.0207)}  # (duration, 2022 dY)

px = pd.concat([load_prices(t, "2021-12-31", "2022-12-31") for t in etfs], axis=1).dropna()
total_ret = px.iloc[-1] / px.iloc[0] - 1

print(f"{'ETF':4} {'duration':>9} {'dY 2022':>8} {'predicted -D*dY':>16} {'actual 2022':>12}")
for t, (d, dy) in etfs.items():
    print(f"{t:4} {d:9.1f} {dy:8.2%} {-d*dy:16.1%} {total_ret[t]:12.1%}")

(px / px.iloc[0] * 100).plot(figsize=(10, 4), color=["steelblue", "darkorange", "crimson"])
plt.title("2022: the duration ladder in real life (indexed to 100)")
plt.ylabel("Value"); plt.tight_layout(); plt.show()
\end{notecode}

\begin{notetable}{@{}p{60bp}p{93bp}p{93bp}p{130bp}p{108.8bp}@{}}
\thcell{ETF} & \thcell{duration} & \thcell{dY 2022} & \thcell{predicted -D*dY} & \thcell{actual 2022} \\
\theadrule
SHY & 1.9 & 3.70\% & -7.0\% & -3.9\% \\ \rowrule
IEF & 7.5 & 2.37\% & -17.8\% & -15.2\% \\ \rowrule
TLT & 17.5 & 2.07\% & -36.2\% & -31.2\% \\
\end{notetable}

\notefigure{figures/bond-duration-convexity/fig2-etf2022.png}
  {Figure 4.3 · 2022 — the duration ladder in real life (indexed to 100)}
  {All three fall through the year, and the drawdown deepens with duration}
  {Source: Author calculations, from the article's computed results.}

The ranking is exactly as duration predicts — and the gaps between predicted and
actual returns are the tutorial's best teaching moment: coupon carry, convexity
(helping TLT), and the fact that ETFs hold rolling portfolios rather than a
single bond all show up in the residual.

\notesubsection{4.4\notenumsep Sanity check: how good is the Taylor approximation?}

Duration-only vs duration+convexity vs exact repricing, for the 30y bond across
shocks up to ±300bp.

\codesize{7}
\begin{notecode}
T = 30
p0 = bond_price(100, 0.04, 0.04, T)
d, c = duration_convexity(100, 0.04, 0.04, T)

shocks = np.linspace(-0.03, 0.03, 61)
exact  = np.array([bond_price(100, 0.04, 0.04 + s, T) / p0 - 1 for s in shocks])
lin    = -d * shocks
quad   = -d * shocks + 0.5 * c * shocks**2

fig, ax = plt.subplots(figsize=(10, 4.5))
ax.plot(shocks * 100, exact * 100, "k",  label="Exact repricing")
ax.plot(shocks * 100, lin   * 100, "--", color="steelblue",  label="Duration only")
ax.plot(shocks * 100, quad  * 100, "--", color="darkorange", label="Duration + convexity")
ax.set_xlabel("Yield shock (%)"); ax.set_ylabel("Price change (%)")
ax.set_title(f"30y bond: Taylor approximation quality (D={d:.1f}, C={c:.0f})")
ax.legend()
plt.tight_layout(); plt.show()

s = 0.02
print(f"+200bp: exact {bond_price(100,0.04,0.06,T)/p0-1:+.2%}   dur-only {-d*s:+.2%}   dur+conv {-d*s+0.5*c*s*s:+.2%}")
\end{notecode}

\notefigure{figures/bond-duration-convexity/fig3-taylor.png}
  {Figure 4.4 · 30y bond — Taylor approximation quality}
  {Duration alone overstates the loss and understates the gain; adding convexity
   closes most of the gap to the exact repricing}
  {Source: Author calculations, from the article's computed results.}

% ========================================================== 5 Conclusion =====
\noteneedroom{330bp}
\notesection{5.\notenumsep Conclusion}

\begin{multicols}{2}
\notesubsection{5a.\notenumsep Strengths}
\begin{notebullets}
\item \textbf{One framework, any bond} — price, duration and convexity summarise
      rate risk in three numbers
\item \textbf{First-order accuracy is excellent} for small yield moves — the
      industry's daily risk language
\item \textbf{Convexity correction} keeps the approximation honest even for
      ±200bp shocks
\item \textbf{Model-free verification} — ETF drawdowns in 2022 line up with
      durations, straight from public data
\end{notebullets}

\notesubsection{5b.\notenumsep Weaknesses \& Limitations}
\begin{notebullets}
\item \textbf{Parallel-shift assumption} — duration assumes the whole curve moves
      together; real curves twist and steepen (key-rate durations fix this)
\item \textbf{Constant yield reinvestment} — YTM assumes coupons reinvest at the
      same rate
\item \textbf{Credit ignored} — this is a rates framework; spreads and default
      risk need their own layer
\item \textbf{ETF $\neq$ single bond} — rolling portfolios have carry and
      roll-down effects the single-bond math misses
\end{notebullets}

\notesubsection{5c.\notenumsep Applications in Practice}
\begin{notebullets}
\item Portfolio duration targeting and immunisation
\item Rate-shock stress tests (the ±100bp tables in every risk report)
\item Barbell vs bullet trades — convexity is the whole game
\end{notebullets}

\notesubsection{5d.\notenumsep Alternatives \& Extensions}
\begin{notebullets}
\item \textbf{Key-rate durations} — sensitivity to individual curve points
\item \textbf{Yield curve bootstrapping with QuantLib} — the natural companion
      piece (curriculum M2)
\item \textbf{Term-structure models} — Vasicek, Hull-White, CIR for simulating
      rates rather than shocking them
\item \textbf{The SVB collapse} — our future case study on what happens when
      duration risk meets deposit flight
\end{notebullets}
\end{multicols}

\end{document}
